\begin{abstract}
We study whether prefix compute allocation in verifiable reasoning can be modeled by ordered confidence thresholding.
To answer that question, we introduce a prefix-level oracle benchmark on two exact-checker domains, arithmetic and linear equations, with three actions: \continue, \revise, and \abstain.
The benchmark reveals substantial high-determinacy action crossing: prefixes with similar continuation scores can still require different oracle-optimal actions.
This is evidence against simple continuation-score thresholding, and more broadly against treating one ordered confidence score as a sufficient description of low-regret control.

We then compare ordered scalar controllers, learned one-dimensional controllers, direct policies, and structured factorized controllers.
The resulting controller picture is conditional rather than universal.
Linear equations strongly disfavors ordered scalar control, while arithmetic is more mixed under repeatability.
For structured controllers, the most stable result is a deployment gap: exact-state factorization is often strong, but deployable predicted-state performance deteriorates sharply under noisy state estimates.
Target-cleaning and exact-state distillation probes recover only domain-specific fractions of this gap and do not restore a general algorithmic win.
That gap appears not only in action regret, but also in revision harm and compute-value calibration.

Our strongest contribution is therefore benchmark- and mechanism-level rather than leaderboard-level:
prefix compute allocation in verifiable reasoning is better understood as \processstate control than as ordered thresholding, exact-checker prefix benchmarks are necessary for testing that claim, and reliable state identification is the main deployment bottleneck for structured controllers.
\end{abstract}
